\documentclass[10pt,a4paper]{amsart}
\usepackage[margin=19mm]{geometry}
\usepackage{amsmath,amssymb,mathtools}
\usepackage[scr=rsfs,cal=euler]{mathalfa}
\usepackage{xfrac}
\usepackage[unicode,hidelinks,bookmarks=false]{hyperref}
\newcommand{\ie}{i.e.\ }
\newcommand{\cf}{cf.\ }
\newcommand{\resp}{respectively\ }
\newcommand{\Subsection}{\subsection}
\providecommand{\nobreakdash}{\nobreak\hbox{-}}
\setlength{\emergencystretch}{2em}
\multlinegap0pt
\usepackage{amssymb}
\usepackage{tikz}
\newtheorem{lemma}{Lemma}
\newtheorem{proposition}{Proposition}
\title{Closing the gap and settling the problem of queens\\ on an $n\times n$ board, each attacking at most one other}
\author{Kristina Ago, Bojan Ba\v{s}i\'c, Radojka Ciganovi\'c}
\date{Source: arXiv:2608.27432v1 (CC BY 4.0)}
\begin{document}
\maketitle

\noindent\emph{Source and licence.} Closing the gap and settling the problem of queens\\ on an $n\times n$ board, each attacking at most one other. Version arXiv:2608.27432v1 (CC BY 4.0), distributed by arXiv under the Creative Commons Attribution 4.0 International licence, http://creativecommons.org/licenses/by/4.0/, as recorded in the arXiv licence metadata for this exact version and shown on the abstract page https://arxiv.org/abs/2608.27432v1. Full licence notice: \texttt{licenses/arxiv-2608.27432-CC-BY-4.0.txt}.\par\smallskip

\noindent\textit{Editorial scope.} Counted unit: the article's proposition split (source0.tex lines 2902--2910). The article's complete original proof of it is delivered with it (source0.tex lines 2912--2929, one \texttt{proof} environment of 18 lines). No mathematics is written, repaired or re-derived: the statement and the proof are the article's own lines, in the article's own notation, and no retained line is substituted or re-flowed.  Retained uncounted as the source-local context the proof needs: lines 2789--2797.  Nothing was omitted from any retained span, so the package carries no omission record.  No retained line contains a citation, so the wrapper carries no bibliography and no bibliography entry.  No retained line contains a figure environment, an image-inclusion command or drawing code, so no image is delivered and the package declares no asset dependency.  Editorial formatting only, in addition: an AMS \texttt{amsart} preamble that restates the article's own declarations and macros used by the retained lines, namely the environments \texttt{lemma}, \texttt{proposition} on separate counters, together with the standard packages those lines need.  The retained environments print in this document as Lemma 1, Proposition 1 in order of appearance, whereas the article prints the same items with the numbering of its own full document.  The complete native source of the article as distributed in the arXiv e-print for this exact version is bundled unchanged as \texttt{sources/208715/source0.tex}.
\begin{lemma}
\label{clearance}
Let $H_A$ be a host and $d$ its clearance. Let $B$ be a positive integer with $\gcd(B,6)=1$. Then
every queen $(r,c)$ of $H_A[B]$ satisfies $|r-c|\geqslant M(d,B)$, where
\[
M(d,B)=dB-\frac{B-1}2.
\]
In other words, the clearance of $H_A[B]$ is at least $M(d,B).$
\end{lemma}

\begin{proposition}
\label{split}
Let $3\mid A$. Let $H$ be an admissible configuration on an $A\times A$ board with $\frac{4A}3$ queens, having no queen on
either of the two main diagonals. Let $B$ be an integer with $B\geqslant5$ and
$\gcd(B,6)=1$, let $t\in\{3, 4, 5\}$, and put $N=AB$. Then the queens
of $H[B]$, shifted $L$ columns
to the right, together with the new queens listed above, form an optimal
configuration on the $(N+t)\times(N+t)$ board.
\end{proposition}

\begin{proof}
The number of queens is right, so only admissibility is at issue. Old and new
queens occupy disjoint sets of rows, and also disjoint sets of columns, so only
differences and sums can cause trouble.

Among the new queens there is nothing to check beyond what the table shows: for
$t=3$ a pair in a column and a pair in a row; for $t=4$ a pair in a row, a pair
in a column, and one queen meeting nobody; for $t=5$ a pair in a row, a pair on
an antidiagonal, and a pair in a column. Apart from these occasions, the four
numbers of the new queens are pairwise different, which one reads off the table
at a glance.

It remains to check that an old and a new queen never attack each other. Clearly, this is never the case for horizontal and vertical attacks, so it remains to take care of diagonal attacks.

Note that all the new queens in the right-hand columns belong to the extension of either the main diagonal of $H[B]$, or a diagonal at distance at most $2$ from the main diagonal. As the clearance of $H$ is at least $1$ (since it was assumed that $H$ has no queens on the main diagonal), Lemma \ref{clearance} gives that the clearance of $H[B]$ is at least $3$. Therefore, we conclude that none of the new queens in the right-hand columns attack an old queen.

Regarding the new queens in the left-hand columns, we note the following. The main antidiagonal of $H[B]$ is empty, as well as the antidiagonal immediately below it. The first claim follows as the main antidiagonal of $H$ is empty (by the assumption), and the second claim follows by the same observation, together with additionally acknowledging that the cell $(1,1)$ is empty in our $B\times B$ board. This completes the proof.
\end{proof}

\end{document}
